\documentclass[11pt]{article}

\usepackage[margin=1in]{geometry}
\usepackage{amsmath,amssymb}
\usepackage{siunitx}
\usepackage{booktabs}
\usepackage{caption}   % for \captionof{table}
\usepackage[hidelinks]{hyperref}

\sisetup{per-mode=symbol}

\title{Hover Lift Estimates for a Two-Blade VTOL Nacelle Rotor}
\author{Jon Farhat (inputs)}
\date{\today}

\begin{document}
\maketitle

\section{Inputs (Per Nacelle)}
\begin{itemize}
  \item Blades: $N_b=2$, variable pitch (0--10$^\circ$)
  \item Chord: $c=\SI{0.030}{m}$
  \item Electrical power in hover: $P_e=\SI{400}{W}$ (about \SI{27}{A} on 4S)
  \item Nacelle mass: $m_n=\SI{1.292}{kg}$, gravity $g=\SI{9.81}{m/s^2}$
  \item Air density: $\rho=\SI{1.225}{kg/m^3}$
  \item Radii evaluated: $R=\SI{0.435}{m}$ and $R=\SI{0.600}{m}$
\end{itemize}

\section{Geometry}
Disk area:
\begin{equation}
A=\pi R^2
\end{equation}
Solidity (shows where chord enters blade-element models):
\begin{equation}
\sigma=\frac{N_b c}{\pi R}
\end{equation}

Computed:
\begin{align}
A_{435} &= \pi(0.435)^2 = \SI{0.594}{m^2}, &
\sigma_{435} &= \frac{2(0.030)}{\pi(0.435)} = 0.0439 \\
A_{600} &= \pi(0.600)^2 = \SI{1.131}{m^2}, &
\sigma_{600} &= \frac{2(0.030)}{\pi(0.600)} = 0.0318
\end{align}

\section{Baseline Weight (Nacelle Only)}
\begin{equation}
W_n=m_n g = 1.292\times 9.81 = \SI{12.67}{N}
\end{equation}

\section{Hover Thrust From Power (Momentum Theory)}
Induced velocity in hover:
\begin{equation}
v_i=\sqrt{\frac{T}{2\rho A}}
\end{equation}
Ideal induced power:
\begin{equation}
P_{i,\mathrm{ideal}} = T v_i = \frac{T^{3/2}}{\sqrt{2\rho A}}
\end{equation}
Solve for thrust vs \emph{shaft} power $P_s$:
\begin{equation}
T_{\mathrm{ideal}}(P_s)=\left(P_s\sqrt{2\rho A}\right)^{2/3}
\end{equation}

We model electrical-to-effective power with an overall effectiveness $\eta$:
\begin{equation}
P_s\approx \eta P_e
\qquad\Rightarrow\qquad
T_{\mathrm{real}}\approx \eta^{2/3} T_{\mathrm{ideal}}(P_e)
\end{equation}

\section{Case Results (Per Nacelle)}
Using $P_e=\SI{400}{W}$:

\subsection*{Case A: $R=\SI{0.435}{m}$}
\begin{align}
T_{\mathrm{ideal},435} &= \left(400\sqrt{2(1.225)(0.594)}\right)^{2/3} = \SI{61.5}{N} \\
\frac{T_{\mathrm{ideal},435}}{g} &= \SI{6.27}{kgf} \\
m_{\mathrm{payload,ideal},435} &= 6.27 - 1.292 = \SI{4.98}{kg}
\end{align}

\subsection*{Case B: $R=\SI{0.600}{m}$}
\begin{align}
T_{\mathrm{ideal},600} &= \left(400\sqrt{2(1.225)(1.131)}\right)^{2/3} = \SI{76.2}{N} \\
\frac{T_{\mathrm{ideal},600}}{g} &= \SI{7.77}{kgf} \\
m_{\mathrm{payload,ideal},600} &= 7.77 - 1.292 = \SI{6.48}{kg}
\end{align}

\section{Realistic Bracket Using \texorpdfstring{$\eta$}{eta}}
Pick example effectiveness values:
\[
\eta\in\{0.35,\,0.50,\,0.60\}
\quad\Rightarrow\quad
\eta^{2/3}\in\{0.497,\,0.630,\,0.711\}.
\]

\bigskip
\noindent\captionof{table}{Hover thrust and payload above nacelle mass (per nacelle) at $P_e=\SI{400}{W}$.}
\begin{center}
\begin{tabular}{@{}lrrrr@{}}
\toprule
Case & $\eta$ & Thrust (kgf) & Payload above nacelle (kg) \\
\midrule
$R=\SI{0.435}{m}$ & Ideal ($\eta=1$) & 6.27 & 4.98 \\
$R=\SI{0.435}{m}$ & 0.35 & 3.11 & 1.82 \\
$R=\SI{0.435}{m}$ & 0.50 & 3.95 & 2.66 \\
$R=\SI{0.435}{m}$ & 0.60 & 4.46 & 3.17 \\
\midrule
$R=\SI{0.600}{m}$ & Ideal ($\eta=1$) & 7.77 & 6.48 \\
$R=\SI{0.600}{m}$ & 0.35 & 3.86 & 2.57 \\
$R=\SI{0.600}{m}$ & 0.50 & 4.90 & 3.60 \\
$R=\SI{0.600}{m}$ & 0.60 & 5.53 & 4.24 \\
\bottomrule
\end{tabular}
\end{center}

\section{Two Nacelles (Total Payload Above Both Nacelles)}
If you have two identical nacelles (two rotors total), approximate total payload above both nacelles:
\[
m_{\mathrm{payload,total}} \approx 2\,m_{\mathrm{payload,per\;nacelle}}
\]
(assuming no large rotor--rotor interference and comparable operating points).

\end{document}
